\section{Implemented Predictor and Conditional Uncertainty}
\label{app:implemented-filter}

This section specifies the local operators and conditional distributions
used by the fixed continuous estimator. It distinguishes identities of
that local model from approximations introduced by its finite history
bank. The numerical checks below exercise these identities on controlled
inputs; they do not recalibrate or rerun the reported evaluation.

\subsection{Affine Operators at a Consistent Topological Anchor}
\label{app:affine-operators}

For network hypothesis $\ell$, let $(\bar x_\ell,\bar u_\ell)$ be the
operating point determined by the prefix-fitted load and dispatch factors.
Fix its differential coordinates while
solving the algebraic constraints separately for each topology $q$:
\begin{equation}
 g_{\ell,q}(\bar x_\ell,\bar z_{\ell,q},\bar u_\ell)=0.
 \label{eq:impl-anchor}
\end{equation}
The resulting point need not be a differential equilibrium after a line
opening. Assume a nonsingular algebraic Jacobian and a fixed angle gauge.
In the expressions below, $f$ is the differential vector field after
division by any nonunit differential mass entries. All derivatives are
evaluated at the topology-specific point in \eqref{eq:impl-anchor}; the
network index is suppressed for readability. Define
\begin{align}
 c_q&=f_q(\bar x,\bar z_q,\bar u),\nonumber\\
 A_q^c&=f_x-f_zg_z^{-1}g_x,&
 B_q^c&=f_u-f_zg_z^{-1}g_u,\nonumber\\
 C_q^c&=H_x-H_zg_z^{-1}g_x,&
 D_q^c&=H_u-H_zg_z^{-1}g_u,\nonumber\\
 d_q^{\rm el}&=H_q(\bar x,\bar z_q,\bar u)
                  -H_0(\bar x,\bar z_0,\bar u).
 \label{eq:impl-affine-continuous}
\end{align}
For the admitted scalar load direction $\delta u=La$, the continuous
approximation is $\dot{\delta x}=A_q^c\delta x+c_q+B_q^cLa$ and
$\delta y=C_q^c\delta x+d_q^{\rm el}+D_q^cLa$.
For a line, the topology selects the operator and the amplitude direction
is zero. A nonzero pre-event state contributes through $\delta x$;
$d_q^{\rm el}$ supplies the zero-deviation offset at the fitted anchor.
This is a local expansion, so large excursions can still cause prediction
error even when the algebraic solve and derivatives are accurate.

For a mode lasting $h$, compute the affine flow without inverting
$A_q^c$:
\begin{align}
 \mathcal A_q&=\begin{pmatrix}A_q^c&c_q&B_q^cL\\0&0&0\\0&0&0\end{pmatrix},\nonumber\\
 \exp(h\mathcal A_q)
 &=\begin{pmatrix}F_q(h)&\beta_q(h)&U_q(h)\\0&1&0\\0&0&1\end{pmatrix}.
 \label{eq:impl-affine-flow}
\end{align}
Thus $\beta_q(h)=\int_0^h\exp(A_q^cs)c_q\,ds$.
If an admitted switch divides one report into an old-mode interval
$h_1$ and a new-mode interval $h_2$, then
\begin{equation}
 F=F_2F_1,\qquad \beta=F_2\beta_1+\beta_2,\qquad
 U=F_2U_1+U_2.
 \label{eq:impl-flow-composition}
\end{equation}
The physical state is continuous across this composition. The final-mode
measurement map uses the consistent algebraic right limit. The predictor
then appends phase history and the deterministic memory of the
instrumental transformation. Scaling these coordinates and applying the
same affine augmentation to \eqref{eq:impl-flow-composition} produces the
stored $T_o,J_o,g_o,u_o,j_o,v_o$ operators. In their report ordering,
\begin{align}
 r_k&=J_o\xi_k+j_o+v_oa+\epsilon_k,\nonumber\\
 \xi_{k+1}&=T_o\xi_k+g_o+u_oa+\eta_k.
 \label{eq:impl-report-order}
\end{align}
The measurement is scored before correcting the state for the next
report. The constant startup discrepancy is applied to the observations;
it does not replace the topology-dependent affine terms.

The construction has three steps. A static solve of
\eqref{eq:impl-anchor} exports the network forcing, the PMU reference and
their reduced Jacobians; the conditional builder forms the topology
offsets, the mode flows and the one-report operators; and a final check
verifies that the constructed reference matches the prefix-fitted
operating point. The reproduction manifest preserves
these files and their hashes. No event truth is used to select the
startup operating point.

\subsection{Four Amplitude Intervals Represent One Prior}
\label{app:severity-partition}

Let $I=[-0.2,0.2]$ be partitioned into four disjoint intervals $I_j$
up to endpoints of zero probability. Write
$Z_J(\mu,v)=\int_J\mathcal N(a;\mu,v)\,da$ for any interval $J$.
The initial prior has the exact decomposition
\begin{align}
 p_0(a)&=\sum_{j=1}^4\alpha_{j,0}
            \operatorname{TN}_{I_j}(a;\mu_0,v_0),\nonumber\\
 \alpha_{j,0}&=Z_{I_j}(\mu_0,v_0)/Z_I(\mu_0,v_0).
 \label{eq:impl-prior-partition}
\end{align}
The implementation assigns one history per source, proposed onset,
restoration path, and interval. Each retains its own
$(m_{h,j},d_{h,j},\mu_{h,j},v_{h,j},w_{h,j})$; there is no per-report
collapse to a single amplitude distribution. The four interval histories
use the same physical operators. They are neither quadrature nodes nor
four separately linearized severity models.

For a fixed physical history with observations $Y$, let
$Z_j(Y)=\int_{I_j}p(Y\mid a,h)
\operatorname{TN}_{I_j}(a;\mu_0,v_0)\,da$. Then
\begin{equation}
 p(Y\mid h)=\sum_j\alpha_{j,0}Z_j(Y),\qquad
 \alpha_j(Y)=\frac{\alpha_{j,0}Z_j(Y)}{p(Y\mid h)}.
 \label{eq:impl-partition-evidence}
\end{equation}
Applying the local predictive integral sequentially gives the same
identity. Identical linear operators and gains preserve the common
untruncated Gaussian parameters for interval copies of the same path;
their truncated moments and normalized weights differ. Subsequent beam
pruning can remove interval histories, so equivalence with an unpartitioned
bank is not guaranteed after pruning. The recorded discarded mass
describes that approximation and is not a bound on all future decisions.

\subsection{Amplitude Dependence and Mixture Compression}
\label{app:conditional-moments}

Condition on a retained history and the reports used so far. For
$a\sim\operatorname{TN}_{[l,u]}(\mu,v)$, set
$\sigma=\sqrt v$, $\alpha=(l-\mu)/\sigma$,
$\beta=(u-\mu)/\sigma$, $Z=\Phi(\beta)-\Phi(\alpha)$, and
$\lambda=[\phi(\alpha)-\phi(\beta)]/Z$. Its moments are
\begin{align}
 \bar a&=\mu+\sigma\lambda,\nonumber\\
 V_a&=v\left[1+\frac{\alpha\phi(\alpha)-\beta\phi(\beta)}{Z}
                         -\lambda^2\right].
 \label{eq:impl-truncated-moments}
\end{align}
The truncated-moment routine evaluates the normalizer in the log
domain and bounds small numerical excursions of the moments.
The Gaussian parameters $\mu,v$ are therefore distinct from the
conditional amplitude mean and variance.

For the affine conditional state model
$\xi\mid a,h\sim\mathcal N(m_h+d_ha,P_h)$, total expectation and
covariance give
\begin{align}
 \bar\xi_h&=m_h+d_h\bar a_h,\nonumber\\
 \operatorname{Cov}(\xi,a\mid h)&=d_hV_{a,h},\nonumber\\
 V_h&=P_h+d_hV_{a,h}d_h^\top.
 \label{eq:impl-state-amplitude-moments}
\end{align}
With normalized history weights $\omega_h$, including the no-new-event
history, the complete mixture moments would be
\begin{align}
 \bar\xi&=\sum_h\omega_h\bar\xi_h,\nonumber\\
 P_{\rm mix}&=\sum_h\omega_h
       \left[V_h+(\bar\xi_h-\bar\xi)(\bar\xi_h-\bar\xi)^\top\right].
 \label{eq:impl-mixture-covariance}
\end{align}
This identity applies within a network hypothesis and its common state
coordinates. It does not justify averaging absolute rotor states between
different network equilibria.

The continuous estimator uses the shared regime covariance $P_\rho$ in
place of $P_h$. It updates the affine means and amplitude distributions,
and computes marginal diagonal state variances for diagnostics. At
closure it carries $\bar\xi$, clears explicit histories, and records the
trace of amplitude/mean spread omitted from the new single-history
representation. It does not carry the complete covariance in
\eqref{eq:impl-mixture-covariance}, nor the uncertainty of the fitted
startup context. Instrumental and discrepancy states remain in the
carried augmented mean.

For comparison, an exact conditional covariance recursion for the local
linear model \eqref{eq:impl-report-order}, with independent noises of
covariances $R,Q$, is
\begin{align}
 S_h&=J_oP_hJ_o^\top+R,\nonumber\\
 K_h&=T_oP_hJ_o^\top S_h^{-1},\nonumber\\
 P_h^+&=T_oP_hT_o^\top+Q-K_hS_hK_h^\top.
 \label{eq:impl-reference-riccati}
\end{align}
This is a diagnostic reference, not a description of a new executed
continuous estimator. The fixed implementation solves the Riccati
equation for the nominal operator in each discrepancy regime and uses
that shared covariance for the admitted switched operators.

Sequential likelihood factorization conditions on past reports;
it does not require marginal independence of the observations. The
instrumental VAR/AR filter and the latent discrepancy AR state already
model temporal dependence. Remaining covariance misspecification,
linearization error, adaptive proposals, and compression can nevertheless
affect the resulting weights. Their empirical calibration requires
pre-update residual and coverage diagnostics, separately from the
algebraic checks below.

\subsection{Deterministic Verification and Its Boundary}
\label{app:math-verification}

A deterministic check loads the implemented amplitude routine without
running a campaign. It compares sequential one-interval
and four-interval evidence and moments with direct scalar quadrature,
including zero information and amplitudes on interval boundaries.
An independent small switched linear system checks affine-flow
composition against numerical integration. A linear Gaussian reference
also compares the conditional affine-state recursion with a joint
state--amplitude Kalman recursion for an untruncated Gaussian amplitude;
the truncated case is checked separately by quadrature. All histories are retained in these
checks. Source hashes, tolerances, and complete results accompany the
generated table.
\begin{table}[t]\centering\small
\caption{Deterministic local-model checks. Maximum absolute errors over 72 scalar updates, five switch positions, and 12 joint Gaussian updates; each row uses its stated reference.}
\label{tab:implemented-filter-checks}
\begin{tabular}{@{}lr@{}}\toprule
Check & Maximum error \\\midrule
Log evidence: partition & $1.78\times10^{-15}$ \\
Amplitude mean/variance & $4.16\times10^{-17}$ \\
Log evidence: quadrature & $3.11\times10^{-15}$ \\
Affine flow: ODE reference & $2.22\times10^{-16}$ \\
Joint mean/covariance & $5.55\times10^{-17}$ \\
\bottomrule\end{tabular}
\end{table}

These checks establish local-model identities on the stated inputs.
They establish neither calibration of the pruned network bank nor
robustness of its nonlinear physical approximation, and add no event
outcomes to the evaluation cohorts.
